\documentclass[margin=5mm]{standalone}
\usepackage[T1]{fontenc}
\usepackage{tikz}
\usepackage{sourcecodepro}
\usetikzlibrary{positioning,calc}
\renewcommand{\familydefault}{\ttdefault}

% Group registry: one definition supplies both the rendering colour and legend.
\newcommand{\PinGroup}[3]{%
  \definecolor{pin-#1}{RGB}{#2}%
  \expandafter\def\csname group-label-#1\endcsname{#3}}
% One definition supplies each group's colour and legend label.
\PinGroup{power}   {239,83,80}   {Power supply}
\PinGroup{ground}  {120,144,156} {Ground}
\PinGroup{analog}  {139,195,74}  {Analog / ADC}
\PinGroup{digital} {66,165,245}  {Digital / GPIO}
\PinGroup{uart}    {255,167,38}  {UART}
\PinGroup{spi}     {171,71,188}  {SPI}
\PinGroup{usb}     {38,198,218}  {USB / USB-C}
\PinGroup{touch}   {0,150,136}   {Capacitive touch}
\PinGroup{control} {141,110,99}  {Reset / interrupt}
\PinGroup{timer}   {190,211,62}  {Timer / capture / PWM}

\definecolor{ink}{RGB}{38,50,56}
\definecolor{muted}{RGB}{100,116,128}
\definecolor{chipfill}{RGB}{247,249,251}

% Rendering stays independent of the physical pin and peripheral data.
\newcommand{\Function}[5]{%
  \ifnum\Side=0
    \node[badge,fill=pin-#1!25,draw=pin-#1!80!ink,
          anchor=east,left=1.2mm of \Previous] (function-\PinNumber-\arabic{functionindex})
          {\if\relax\detokenize{#2}\relax\else #2\,\fi#3#4%
           \if\relax\detokenize{#5}\relax\else\;{\tiny[#5]}\fi};%
  \else
    \node[badge,fill=pin-#1!25,draw=pin-#1!80!ink,
          anchor=west,right=1.2mm of \Previous] (function-\PinNumber-\arabic{functionindex})
          {\if\relax\detokenize{#2}\relax\else #2\,\fi#3#4%
           \if\relax\detokenize{#5}\relax\else\;{\tiny[#5]}\fi};%
  \fi
  \edef\Previous{function-\PinNumber-\arabic{functionindex}}%
  \stepcounter{functionindex}}
\newcounter{functionindex}

\newcommand{\Pin}[4]{%
  \def\PinNumber{#1}%
  \ifnum#1<9
    \def\Side{0}\pgfmathsetmacro{\PinY}{8-#1}%
    \path[draw=ink,fill=pin-#3] (-0.4,\PinY-0.25) rectangle (0,\PinY+0.25);
    \node[font=\scriptsize,text=muted] at (0.3,\PinY) {#1};
    \node[pinname,fill=pin-#3!20,anchor=east] (pin-#1) at (-0.58,\PinY) {#2};%
  \else
    \def\Side{1}\pgfmathsetmacro{\PinY}{#1-9}%
    \path[draw=ink,fill=pin-#3] (2.6,\PinY-0.25) rectangle (3,\PinY+0.25);
    \node[font=\scriptsize,text=muted] at (2.3,\PinY) {#1};
    \node[pinname,fill=pin-#3!20,anchor=west] (pin-#1) at (3.18,\PinY) {#2};%
  \fi
  \edef\Previous{pin-#1}\setcounter{functionindex}{0}%
  #4}

\begin{document}
\begin{tikzpicture}[
  x=1cm,y=0.8cm,text=ink,
  pinname/.style={draw=ink!65,line width=0.4pt,rounded corners=1pt,
    minimum width=11mm,minimum height=5.5mm,font=\small\bfseries},
  badge/.style={line width=0.35pt,rounded corners=1pt,inner xsep=3pt,
    inner ysep=3pt,minimum height=5.5mm,font=\scriptsize}]
\pagecolor{white}
\path[fill=chipfill,draw=ink,line width=0.8pt,rounded corners=2pt]
  (0,-0.7) rectangle (2.6,7.7);
% Pin 1 marker and package notch.
\fill[ink] (0.3,7.45) circle (0.055);
\draw[ink,line width=0.8pt] (0.95,7.7) arc[start angle=180,end angle=360,x radius=0.35,y radius=0.35];
\node[rotate=90,font=\Large\bfseries,text=ink!75] at (1.3,3.5) {CH552G};
% CH552G, SOP-16, top view. WCH CH552 datasheet V1G, sections 3, 4, 10.4.
% \Pin{physical number}{name}{base group}{functions}
% \Function{group}{peripheral instance}{signal}{channel}{route}
% D/A = default/alternate mapping in PIN_FUNC; empty = fixed pin function.
% The datasheet's underscore suffix means alternate mapping, not active-low.
% All Pxn names support GPIO; shared functions require peripheral configuration.

\Pin{1}{P3.2}{analog}{%
  \Function{analog}{ADC}{AIN}{3}{}
  \Function{uart}{UART1}{TXD}{}{A}
  \Function{control}{}{INT}{0}{}
  \Function{usb}{USB-C}{VBUS}{1}{}}
\Pin{2}{P1.4}{analog}{%
  \Function{analog}{ADC}{AIN}{1}{}
  \Function{spi}{SPI0}{SCS}{}{}
  \Function{timer}{T2}{CLK}{}{A}
  \Function{timer}{T2}{CAP}{1}{A}
  \Function{touch}{TKEY}{TIN}{2}{}
  \Function{usb}{USB-C}{UCC}{1}{}}
\Pin{3}{P1.5}{analog}{%
  \Function{analog}{ADC}{AIN}{2}{}
  \Function{spi}{SPI0}{MOSI}{}{}
  \Function{timer}{}{PWM}{1}{D}
  \Function{touch}{TKEY}{TIN}{3}{}
  \Function{usb}{USB-C}{UCC}{2}{}}
\Pin{4}{P1.6}{digital}{%
  \Function{spi}{SPI0}{MISO}{}{}
  \Function{uart}{UART1}{RXD}{}{D}
  \Function{touch}{TKEY}{TIN}{4}{}}
\Pin{5}{P1.7}{digital}{%
  \Function{spi}{SPI0}{SCK}{}{}
  \Function{uart}{UART1}{TXD}{}{D}
  \Function{touch}{TKEY}{TIN}{5}{}}
\Pin{6}{RST}{control}{%
  \Function{control}{}{RESET}{}{}
  \Function{timer}{T2}{EX}{}{A}
  \Function{timer}{T2}{CAP}{2}{A}}
\Pin{7}{P3.1}{digital}{%
  \Function{uart}{UART0}{TXD}{}{D}
  \Function{timer}{}{PWM}{2}{A}}
\Pin{8}{P3.0}{digital}{%
  \Function{uart}{UART0}{RXD}{}{D}
  \Function{timer}{}{PWM}{1}{A}}
\Pin{9}{P1.1}{analog}{%
  \Function{analog}{ADC}{AIN}{0}{}
  \Function{timer}{T2}{EX}{}{D}
  \Function{timer}{T2}{CAP}{2}{D}
  \Function{touch}{TKEY}{TIN}{1}{}
  \Function{usb}{USB-C}{VBUS}{2}{}}
\Pin{10}{P3.3}{digital}{%
  \Function{control}{}{INT}{1}{}}
\Pin{11}{P3.4}{digital}{%
  \Function{timer}{}{PWM}{2}{D}
  \Function{uart}{UART1}{RXD}{}{A}
  \Function{timer}{T0}{IN}{}{}}
\Pin{12}{P3.6}{usb}{%
  \Function{usb}{USB}{D+ (UDP)}{}{}}
\Pin{13}{P3.7}{usb}{%
  \Function{usb}{USB}{D- (UDM)}{}{}}
\Pin{14}{GND}{ground}{%
  \Function{ground}{}{VSS}{}{}}
\Pin{15}{VCC}{power}{%
  \Function{power}{}{VDD}{}{}}
\Pin{16}{V33}{power}{%
  \Function{power}{USB}{3.3V}{}{}}


% Use the actual badge extents to centre the heading and legend.
\coordinate (bottomleft) at (current bounding box.south west);
\coordinate (bottomright) at (current bounding box.south east);
\coordinate (heading) at ($(current bounding box.north west)!0.5!(current bounding box.north east)$);
\node[anchor=south,font=\Large\bfseries] at ($(heading)+(0,0.65)$) {CH552G / PINOUT};
\node[anchor=south,font=\small,text=muted] at ($(heading)+(0,0.12)$)
  {16-pin SOP \quad | \quad top view};

\newlength{\DiagramWidth}
\path let \p1=(bottomleft),\p2=(bottomright) in
  \pgfextra{\global\DiagramWidth=\dimexpr\x2-\x1\relax};
\coordinate (legendstart) at ($(bottomleft)+(0,-0.7)$);
\foreach \key [count=\i from 0] in {power,ground,analog,digital,uart,spi,usb,touch,control,timer}{%
  \pgfmathtruncatemacro{\col}{mod(\i,5)}%
  \pgfmathtruncatemacro{\row}{int(\i/5)}%
  \pgfmathsetmacro{\fraction}{\col/5}%
  \coordinate (legendcol) at ($(bottomleft)!\fraction!(bottomright)$);
  \node[badge,anchor=north west,fill=pin-\key!25,draw=pin-\key!80!ink]
    at ($(legendcol)+(0,-0.7-\row*0.85)$) {\csname group-label-\key\endcsname};}
\node[anchor=north west,align=left,font=\scriptsize,text=muted,text width=\DiagramWidth]
  at ($(legendstart)+(0,-1.8)$) {%
  D / A = default / alternate mapping in PIN\_FUNC. Alternate labels replace the datasheet's underscore suffix.\\[2pt]
  Pxn names support GPIO; USB data pins require USB\_IO\_EN = 0 for GPIO. Shared labels are selectable functions.\\[2pt]
  T2 CLK = T2 counter input / clock output; T2 EX = T2EX reload/capture input; SCS = SPI slave-select input.\\[2pt]
  USB D+/D- also carry ISP bootloader traffic. P3.6/P3.7 use V33 levels; V33 is not a general 5 V supply pin.\\[2pt]
  SOP-16 does not expose P1.0, P1.2, P1.3 or P3.5: no external crystal pins or T1 input; T2 CLK/CAP1 needs route A.\\[2pt]
  No hardware I2C peripheral. Source: WCH CH552 datasheet V1G, sections 3, 4 and 10.4.};
\end{tikzpicture}
\end{document}
